% Part: propositional-logic

\documentclass[../../include/open-logic-part]{subfiles}

\begin{document}

\olpart{pl}{Logique propositionnelle}

\begin{editorial}
Cette partie porte sur la logique propositionnelle classique. Le premier
chapitre est assez sommaire : il présente essentiellement des définitions
et des résultats, dont beaucoup de démonstrations sont laissées en
exercice. Les développements sur les systèmes de preuve et le théorème de
complétude proviennent de la partie consacrée à la logique du premier
ordre, avec la balise \texttt{FOL} désactivée. Cela supprime tout ce qui
concerne les prédicats, les termes et les quantificateurs, et remplace les
!!{structure}s~$\Struct{M}$ par des !!{valuation}s~$\pAssign{v}$.

Il est prévu d'étoffer cette partie et d'y ajouter d'autres sujets et
résultats, notamment la complétude fonctionnelle des connecteurs.
\end{editorial}

\olimport[syntax-and-semantics]{syntax-and-semantics}

\tagfalse{FOL}

\olimport[../first-order-logic/proof-systems]{proof-systems}

\iftag{prfSC}{%
  \olimport[../first-order-logic/sequent-calculus]{sequent-calculus}
}{}

\iftag{prfND}{%
  \olimport[../first-order-logic/natural-deduction]{natural-deduction}
}{}

\iftag{prfTab}{%
  \olimport[../first-order-logic/tableaux]{tableaux}
}{}

\iftag{prfAX}{%
  \olimport[../first-order-logic/axiomatic-deduction]{axiomatic-deduction}
}{}

\olimport[../first-order-logic/completeness]{completeness}

\tagtrue{FOL}

\OLEndPartHook

\end{document}
